\documentclass[12pt]{article}
\usepackage[utf8]{inputenc}
\usepackage[T1]{fontenc}
\usepackage{amsmath, amssymb, geometry, booktabs}
\usepackage{enumitem}
\geometry{margin=1in}
\setlength{\parindent}{0pt}
\renewcommand{\arraystretch}{1.2}


\title{Solutions -- Tutorial 2 --The Instruments of Trade Policy}
\author{ECON 308 -- International Economic Relations}
\date{Fall 2025}

\begin{document}
\maketitle

\section*{Problem 1 – Tariff and Welfare for a Small Country}

\textbf{Given:}
\[
Q_S = P - 100,\qquad Q_D = 700 - P,\qquad P_W = 200.
\]
Small country $\Rightarrow$ world price fixed at \$200; tariff does not affect $P_W$.

\subsection*{(1) Free-trade equilibrium imports}

Under free trade, domestic price equals world price:
\[
P^{FT} = P_W = 200.
\]
Domestic supply:
\[
Q_S^{FT} = 200 - 100 = 100.
\]
Domestic demand:
\[
Q_D^{FT} = 700 - 200 = 500.
\]
Imports:
\[
M^{FT} = Q_D^{FT} - Q_S^{FT} = 500 - 100 = 400.
\]
\[
\boxed{M^{FT} = 400 \text{ tons.}}
\]

\subsection*{(2) Equilibrium with a \$20 specific tariff}

Tariff $t = 20$ per ton. For a small country, world price remains \$200; domestic consumers pay:
\[
P^{T} = P_W + t = 200 + 20 = 220.
\]
Supply at $P^T$:
\[
Q_S^{T} = 220 - 100 = 120.
\]
Demand at $P^T$:
\[
Q_D^{T} = 700 - 220 = 480.
\]
Imports:
\[
M^{T} = Q_D^{T} - Q_S^{T} = 480 - 120 = 360.
\]
\[
\boxed{P^T = 220,\; Q_S^{T} = 120,\; Q_D^{T} = 480,\; M^{T} = 360.}
\]

\subsection*{(3) Tariff revenue}

Government revenue:
\[
\text{Rev} = t \cdot M^{T} = 20 \times 360 = 7{,}200.
\]
\[
\boxed{\text{Tariff revenue} = \$7{,}200.}
\]

\subsection*{(4) Deadweight losses from production and consumption}

The tariff wedges domestic price $P^T$ above $P_W$ by $t=20$.  
Quantity changes:
\[
\Delta Q_S = Q_S^{T} - Q_S^{FT} = 120 - 100 = 20,
\]
\[
\Delta Q_D = Q_D^{FT} - Q_D^{T} = 500 - 480 = 20.
\]

\textbf{Production distortion (over-production) loss:}
Triangle at the left side of the import demand diagram with base $\Delta Q_S$ and height $t$:
\[
\text{DWL}_{\text{prod}} = \frac{1}{2} \times t \times \Delta Q_S
= \frac{1}{2} \times 20 \times 20 = 200.
\]

\textbf{Consumption distortion (under-consumption) loss:}
Triangle at the right side with base $\Delta Q_D$ and height $t$:
\[
\text{DWL}_{\text{cons}} = \frac{1}{2} \times t \times \Delta Q_D
= \frac{1}{2} \times 20 \times 20 = 200.
\]

\[
\boxed{\text{DWL}_{\text{prod}} = 200,\quad \text{DWL}_{\text{cons}} = 200.}
\]

\subsection*{(5) Net welfare loss}

Total deadweight loss:
\[
\text{Total DWL} = \text{DWL}_{\text{prod}} + \text{DWL}_{\text{cons}} = 200 + 200 = 400.
\]
In a small country, there is \emph{no terms-of-trade gain}, so the tariff simply transfers surplus from consumers to producers and government, plus this net loss:
\[
\boxed{\text{Net welfare loss} = \$400.}
\]

\newpage

\section*{Problem 2 – Large Country Tariff and Terms of Trade}

\textbf{Given:}
\[
P^* = 100 + 0.5Q^* \quad \text{(foreign export supply)}, \qquad Q_M^D = 400 - 2P \quad \text{(Home import demand)}.
\]
$P^*$ is the price received by foreign exporters (world price).  
$P$ is the price paid by Home importers (domestic import price).

\subsection*{(1) Free-trade equilibrium price and imports}

Under free trade, there is no tariff, so:
\[
P = P^* = P^{FT}, \qquad Q = Q^* = Q_M^D.
\]

Use Home import demand:
\[
Q = 400 - 2P.
\]
Use foreign export supply:
\[
P = 100 + 0.5Q.
\]

Substitute $Q = 400 - 2P$ into the export supply equation:
\[
P = 100 + 0.5(400 - 2P) = 100 + 200 - P = 300 - P.
\]
So
\[
2P = 300 \Rightarrow P^{FT} = 150.
\]
Then
\[
Q^{FT} = 400 - 2(150) = 400 - 300 = 100.
\]

\[
\boxed{P^{FT} = 150,\quad Q^{FT} = 100 \text{ units of imports.}}
\]

\subsection*{(2) \$20 tariff: new world price, domestic price, and imports}

Let the tariff be $t = 20$. Then
\[
P = P^* + t.
\]
Equilibrium conditions:
\[
Q = Q_M^D = 400 - 2P, 
\quad
P^* = 100 + 0.5Q, 
\quad
P = P^* + 20.
\]

Plug $P = P^* + 20$ into $Q_M^D$:
\[
Q = 400 - 2(P^* + 20) = 400 - 2P^* - 40 = 360 - 2P^*.
\]
Foreign export supply:
\[
P^* = 100 + 0.5Q = 100 + 0.5(360 - 2P^*) = 100 + 180 - P^* = 280 - P^*.
\]
So
\[
2P^* = 280 \Rightarrow P^* = 140.
\]
Then domestic price:
\[
P = P^* + t = 140 + 20 = 160.
\]
Import quantity:
\[
Q = 360 - 2P^* = 360 - 2(140) = 360 - 280 = 80.
\]

\[
\boxed{P^* = 140,\quad P = 160,\quad Q^{T} = 80.}
\]

Note: World price falls from 150 to 140; Home improves its terms of trade.

\subsection*{(3) Tariff revenue and terms-of-trade gain}

\textbf{Tariff revenue:}
\[
\text{Rev} = t \cdot Q^{T} = 20 \times 80 = 1{,}600.
\]
\[
\boxed{\text{Tariff revenue} = \$1{,}600.}
\]

\textbf{Terms-of-trade (ToT) gain:}  
Before tariff, foreign received $P^{FT}=150$; after, $P^* = 140$.  
The drop in the price Home pays to foreigners (in world markets) is
\[
\Delta P_{\text{ToT}} = 150 - 140 = 10.
\]
Home still imports $Q^T = 80$ units at this lower foreign price, so the ToT gain is:
\[
\text{ToT gain} = \Delta P_{\text{ToT}} \cdot Q^{T} = 10 \times 80 = 800.
\]
\[
\boxed{\text{Terms-of-trade gain} = \$800.}
\]

\subsection*{(4) Welfare comparison: ToT gain vs. deadweight loss}

The total welfare effect of the tariff for a large country is:
\[
\Delta W = \text{ToT gain} - \text{DWL}.
\]

The tariff reduces imports from $Q^{FT}=100$ to $Q^T=80$, a change of:
\[
\Delta Q = 100 - 80 = 20.
\]
Tariff wedge is $t = 20$. In the standard trade diagram (quantity on horizontal, prices on vertical), two small triangles arise at the two edges of the new trade volume. Each has base $\Delta Q$ and height $t$:
\[
\text{DWL triangle (each)} = \frac{1}{2} \times t \times \Delta Q
= \frac{1}{2} \times 20 \times 20 = 200.
\]
Total deadweight loss:
\[
\text{DWL}_{\text{total}} = 2 \times 200 = 400.
\]

Thus,
\[
\Delta W = 800 - 400 = 400 > 0.
\]
\[
\boxed{\text{Net welfare gain} = \$400 \text{ for Home.}}
\]

\subsection*{(5) Why retaliation limits optimal tariffs}

In reality, trading partners typically retaliate by imposing their own tariffs or trade barriers. This:
\begin{itemize}
\item Reduces demand for Home exports,
\item Lowers Home’s export prices,
\item Shrinks or reverses the terms-of-trade gain.
\end{itemize}
Hence the net effect can easily become negative for both countries, making aggressive use of ``optimal tariffs'' politically and economically risky. This is a key reason behind multilateral agreements (e.g., WTO) that discourage such behaviour.

\bigskip
\hrule
\vspace{0.8em}

\section*{Problem 3 – Effective Rate of Protection}

\textbf{Given:}
\begin{itemize}
\item World price of car: \$20{,}000.
\item World price of imported engine: \$8{,}000.
\item Tariff on engine: 25\%.
\item Tariff on finished car: 10\%.
\end{itemize}

Let $V_A$ be domestic value added before protection and $V_A'$ after protection.

\subsection*{(1) Nominal protection rate}

The nominal tariff on the finished car is 10\%.  
\[
\boxed{\text{Nominal protection rate on the car} = 10\%.}
\]

\subsection*{(2) Effective Rate of Protection (ERP)}

\textbf{Without tariffs (free trade):}
\[
P_{\text{car}}^{FT} = 20{,}000,\quad P_{\text{engine}}^{FT} = 8{,}000.
\]
Domestic value added is:
\[
V_A = P_{\text{car}}^{FT} - P_{\text{engine}}^{FT} = 20{,}000 - 8{,}000 = 12{,}000.
\]

\textbf{With tariffs:}

Engine tariff 25\%:
\[
P_{\text{engine}}^{T} = 8{,}000 (1 + 0.25) = 10{,}000.
\]

Car tariff 10\%:
\[
P_{\text{car}}^{T} = 20{,}000 (1 + 0.10) = 22{,}000.
\]

New domestic value added:
\[
V_A' = P_{\text{car}}^{T} - P_{\text{engine}}^{T} = 22{,}000 - 10{,}000 = 12{,}000.
\]

\textbf{Compute ERP:}
\[
\text{ERP} = \frac{V_A' - V_A}{V_A} 
= \frac{12{,}000 - 12{,}000}{12{,}000}= 0.
\]

\[
\boxed{\text{Effective rate of protection} = 0\%.}
\]

\subsection*{(3) Interpretation}

Even though there is a 10\% \emph{nominal} tariff on the final car, the protection to \emph{domestic value added} is \textbf{zero}:
\begin{itemize}
\item The tariff on the input (engine) raises its cost to domestic producers.
\item The tariff on the final good raises the car price.
\item In this numerical example, these effects cancel out so that domestic value added is unchanged at \$12,000.
\end{itemize}
Thus the industry is \emph{not} effectively protected, despite the nominal 10\% tariff. This illustrates why policymakers must look at the \textbf{ERP}, not just nominal tariffs, to understand true protection of domestic production.

\end{document}
